\section{Full Derivations}
\label{app:deriv}

\subsection{Max-pooling Jacobian}
\label{app:maxpool}
Let a pooling window hold inputs $x_1,\dots,x_n$ (row-major; for the brief
$n=4$, $x=(8,7,12,9)$). The forward map is $y=\max_k x_k$. Assume a unique
maximiser $k^\star=\arg\max_k x_k$ (here $k^\star=3$, $x_{k^\star}=12$). For
$\varepsilon$ small enough that the ordering is preserved:
\begin{align}
  y(x_{k^\star}+\varepsilon) &= x_{k^\star}+\varepsilon
    &&\Rightarrow\quad \frac{\partial y}{\partial x_{k^\star}} = 1, \\
  y(x_{k}+\varepsilon) &= x_{k^\star} \ \ (k\neq k^\star)
    &&\Rightarrow\quad \frac{\partial y}{\partial x_{k}} = 0 .
\end{align}
Hence, reshaped to the $2\times2$ block,
\begin{equation}
  \frac{\partial\,\text{Output}}{\partial\,\text{Input}}
  = \begin{bmatrix} 0 & 0 \\ 1 & 0 \end{bmatrix}.
\end{equation}
In back-propagation the upstream gradient $g$ of the pooled cell is therefore
\emph{routed} entirely to position $k^\star$ and every other input receives
zero: max-pooling is a data-dependent hard router / gate. (Ties are resolved to
the first maximiser; the sub-gradient is any convex combination summing to
$g$.) By contrast, for $y=\frac1n\sum_k x_k$ we get
$\partial y/\partial x_k = 1/n = 0.25\ \forall k$, so average pooling
distributes the gradient uniformly and every input contributes to learning.

\subsection{Scale-space causality via the heat equation}
\label{app:causality}
Let $u(x,t) = (G_{\sqrt{t}} * f)(x)$ be the Gaussian scale-space of a 1-D
signal $f$. It satisfies the 1-D heat equation
\begin{equation}
  \frac{\partial u}{\partial t} = \frac12 \frac{\partial^2 u}{\partial x^2},
  \qquad u(x,0)=f(x),
\end{equation}
with $t=\sigma^2$. Differentiating $k$ times in $x$ and writing
$v=\partial_x^k u$ shows $v$ obeys the same equation. The heat equation
satisfies a \emph{maximum principle}: an interior local maximum of $v$ cannot
increase and an interior local minimum cannot decrease with $t$, because at a
local max $\partial_{xx} v \le 0 \Rightarrow \partial_t v \le 0$. Consequently
the number of local extrema of $v$ is non-increasing in $t$
(Hummel--Gidas, Lindeberg~\cite{lindeberg1994}).

Edges are the zero-crossings of $u_{xx}$. A new zero-crossing of $u_{xx}$ could
only appear if $u_{xx}$ developed a new extremum that then crossed zero --- but
$u_{xxx}$'s extrema (hence $u_{xx}$'s monotone segments) cannot increase in
number. Therefore no new zero-crossing (edge) is created as $\sigma$ grows;
existing zero-crossings move continuously and vanish in pairs when two
neighbouring extrema of $u_{xx}$ merge. This is exactly the behaviour observed
in \cref{fig:causality}: the fingerprint is a family of arcs that only close
upward, and the count is monotone non-increasing. It justifies coarse-to-fine
edge tracking: every coarse-scale edge has a unique fine-scale ancestor.

\emph{Caveat.} Strict monotonicity is a 1-D result. In 2-D the isotropic
Gaussian is still the unique separable, rotationally-symmetric kernel that
creates no new structure, but pathological 2-D counter-examples to strict
extremum monotonicity exist; the edge-\emph{count} trend remains
overwhelmingly decreasing, which our measurement (a $1.15\times$ slack in the
monotonicity check) reflects.
